% GENERATED FILE. Edit canonical lessons, then run npm run generate:books.
% canonical-source-hash: fnv1a64:9bd9c9ca37e74b0e
% canonical-lessons: BN-C480-pushti, BN-C480-dhen, BN-C480-nesha, BN-C480-dhumpan, BN-C480-jhorano

\chapter{Staying Healthy}
\label{ch:bn-staying-healthy}
\hlchaptermodality{480}

\begin{chapteropening}
\textbf{By the end of this chapter:} \emph{I can talk about nutrition, meditation and bad habits.}

Complete the last lesson of chapter 480: I can talk about nutrition, meditation and bad habits.
\end{chapteropening}

\section[\texorpdfstring{\bn{পুষ্টি} (pushṭi)}{পুষ্টি (pushṭi)}]{\bn{পুষ্টি} (pushṭi) — nutrition; nourishment}

\label{lesson:BN-C480-pushti}



\begin{quote}
Before the new one: say the Bengali for to nod off, then the Bengali for to stretch on waking.
\end{quote}

\subsection*{You'll want to know: \bn{পুষ্টি}}
\textbf{\bn{পুষ্টি}} — \emph{pushṭi} — \textquotedblleft{}nutrition; nourishment\textquotedblright{}.

\begin{quote}
\bn{শিশুদের} \bn{পুষ্টির} \bn{জন্য} \bn{দুধ} \bn{আর} \bn{ডিম} \bn{দরকার।}
\end{quote}

\begin{quote}
\emph{shishuder pushṭir jônyô dudh ār ḍim dôrkār.}
\end{quote}

\begin{quote}
Children need milk and eggs for nourishment.
\end{quote}

\begin{cousinweb}
\textbf{\bn{পুষ্টি}} (\emph{pushṭi}) is a tatsama, Sanskrit \emph{puṣṭi}, \textquotedblleft{}a well-nourished condition, growth, thriving\textquotedblright{} (Monier-Williams).
\end{cousinweb}

\begin{grammarlens}[title={What you've built}]
One more word for describing what you lived and felt.
\end{grammarlens}

\subsection*{Your turn}
\begin{itemize}
\raggedright
  \item \emph{Say it:} \emph{pushṭi}
  \item \emph{Say it:} \emph{pushṭi}, once more
  \item \emph{Say it:} \emph{shishuder pushṭir jônyô dudh ār ḍim dôrkār.} — Children need milk and eggs for nourishment.
  \item \emph{Recall:} say the Bengali for to nod off, then the Bengali for to stretch on waking, then say \emph{pushṭi} again
\end{itemize}

\begin{tcolorbox}[breakable,colback=teal!4,colframe=teal!35!black,title={Before you move on}]
What does \bn{পুষ্টি} mean? (\textquotedblleft{}Nutrition; nourishment\textquotedblright{}.) Say it once more.
\end{tcolorbox}
\section[\texorpdfstring{\bn{ধ্যান} (dhæn)}{ধ্যান (dhæn)}]{\bn{ধ্যান} (dhæn) — meditation}

\label{lesson:BN-C480-dhen}



\begin{quote}
Before the new one: say the Bengali for to stretch on waking, then the Bengali for nutrition.
\end{quote}

\subsection*{You'll want to know: \bn{ধ্যান}}
\textbf{\bn{ধ্যান}} — \emph{dhæn} — \textquotedblleft{}meditation\textquotedblright{}.

\begin{quote}
\bn{দাদু} \bn{রোজ} \bn{ভোরে} \bn{ধ্যান} \bn{করেন।}
\end{quote}

\begin{quote}
\emph{dādu roj bhore dhæn kôren.}
\end{quote}

\begin{quote}
Grandfather meditates every day at dawn.
\end{quote}

\begin{cousinweb}
\textbf{\bn{ধ্যান}} (\emph{dhæn}) is a tatsama, Sanskrit \emph{dhyāna}, \textquotedblleft{}meditation, reflection\textquotedblright{} (Monier-Williams).
\end{cousinweb}

\begin{grammarlens}[title={What you've built}]
One more word for describing what you lived and felt.
\end{grammarlens}

\subsection*{Your turn}
\begin{itemize}
\raggedright
  \item \emph{Say it:} \emph{dhæn}
  \item \emph{Say it:} \emph{dhæn}, once more
  \item \emph{Say it:} \emph{dādu roj bhore dhæn kôren.} — Grandfather meditates every day at dawn.
  \item \emph{Recall:} say the Bengali for to stretch on waking, then the Bengali for nutrition, then say \emph{dhæn} again
\end{itemize}

\begin{tcolorbox}[breakable,colback=teal!4,colframe=teal!35!black,title={Before you move on}]
What does \bn{ধ্যান} mean? (\textquotedblleft{}Meditation\textquotedblright{}.) Say it once more.
\end{tcolorbox}
\section[\texorpdfstring{\bn{নেশা} (neshā)}{নেশা (neshā)}]{\bn{নেশা} (neshā) — an addiction; intoxication}

\label{lesson:BN-C480-nesha}



\begin{quote}
Before the new one: say the Bengali for nutrition, then the Bengali for meditation.
\end{quote}

\subsection*{You'll want to know: \bn{নেশা}}
\textbf{\bn{নেশা}} — \emph{neshā} — \textquotedblleft{}an addiction; intoxication\textquotedblright{}.

\begin{quote}
\bn{ফোন} \bn{দেখা} \bn{অরুণের} \bn{একটা} \bn{নেশা।}
\end{quote}

\begin{quote}
\emph{phon dækhā ôruner ekṭā neshā.}
\end{quote}

\begin{quote}
Looking at his phone is an addiction for Arun.
\end{quote}

\begin{cousinweb}
\textbf{\bn{নেশা}} (\emph{neshā}) is usually traced to Persian \emph{nasha}, \textquotedblleft{}intoxication\textquotedblright{}.
\end{cousinweb}

\begin{grammarlens}[title={What you've built}]
One more word for describing what you lived and felt.
\end{grammarlens}

\subsection*{Your turn}
\begin{itemize}
\raggedright
  \item \emph{Say it:} \emph{neshā}
  \item \emph{Say it:} \emph{neshā}, once more
  \item \emph{Say it:} \emph{phon dækhā ôruner ekṭā neshā.} — Looking at his phone is an addiction for Arun.
  \item \emph{Recall:} say the Bengali for nutrition, then the Bengali for meditation, then say \emph{neshā} again
\end{itemize}

\begin{tcolorbox}[breakable,colback=teal!4,colframe=teal!35!black,title={Before you move on}]
What does \bn{নেশা} mean? (\textquotedblleft{}An addiction; intoxication\textquotedblright{}.) Say it once more.
\end{tcolorbox}
\section[\texorpdfstring{\bn{ধূমপান} (dhūmpān)}{ধূমপান (dhūmpān)}]{\bn{ধূমপান} (dhūmpān) — smoking (tobacco)}

\label{lesson:BN-C480-dhumpan}



\begin{quote}
Before the new one: say the Bengali for meditation, then the Bengali for an addiction.
\end{quote}

\subsection*{You'll want to know: \bn{ধূমপান}}
\textbf{\bn{ধূমপান}} — \emph{dhūmpān} — \textquotedblleft{}smoking (tobacco)\textquotedblright{}.

\begin{quote}
\bn{ধূমপান} \bn{স্বাস্থ্যের} \bn{জন্য} \bn{খারাপ।}
\end{quote}

\begin{quote}
\emph{dhūmpān svāsthyer jônyô khārāp.}
\end{quote}

\begin{quote}
Smoking is bad for your health.
\end{quote}

\bn{ধূমপান} \bn{করা}, to smoke.

\begin{cousinweb}
\textbf{\bn{ধূমপান}} (\emph{dhūmpān}) is a tatsama, Sanskrit \emph{dhūmapāna}, \textquotedblleft{}drinking smoke\textquotedblright{}, inhaling smoke and, later, smoking tobacco (Monier-Williams).
\end{cousinweb}

\begin{grammarlens}[title={What you've built}]
One more word for describing what you lived and felt.
\end{grammarlens}

\subsection*{Your turn}
\begin{itemize}
\raggedright
  \item \emph{Say it:} \emph{dhūmpān}
  \item \emph{Say it:} \emph{dhūmpān}, once more
  \item \emph{Say it:} \emph{dhūmpān svāsthyer jônyô khārāp.} — Smoking is bad for your health.
  \item \emph{Recall:} say the Bengali for meditation, then the Bengali for an addiction, then say \emph{dhūmpān} again
\end{itemize}

\begin{tcolorbox}[breakable,colback=teal!4,colframe=teal!35!black,title={Before you move on}]
What does \bn{ধূমপান} mean? (\textquotedblleft{}Smoking (tobacco)\textquotedblright{}.) Say it once more.
\end{tcolorbox}
\section[\texorpdfstring{\bn{ঝরানো} (jhôrāno)}{ঝরানো (jhôrāno)}]{\bn{ঝরানো} (jhôrāno) — to shed (weight, tears)}

\label{lesson:BN-C480-jhorano}



\begin{quote}
Before the new one: say the Bengali for an addiction, then the Bengali for smoking.
\end{quote}

\subsection*{You'll want to know: \bn{ঝরানো}}
\textbf{\bn{ঝরানো}} — \emph{jhôrāno} — \textquotedblleft{}to shed (weight, tears)\textquotedblright{}.

\begin{quote}
\bn{ব্যায়াম} \bn{করে} \bn{বাবা} \bn{পাঁচ} \bn{কেজি} \bn{ওজন} \bn{ঝরালেন।}
\end{quote}

\begin{quote}
\emph{bæyām kôre bābā pā̃ch keji ojôn jhôrālen.}
\end{quote}

\begin{quote}
By exercising, my father shed five kilos.
\end{quote}

Past (he/she …): \bn{তিনি} \bn{দশ} \bn{কেজি} \bn{ওজন} \bn{ঝরালেন}, \emph{tini dôsh keji ojôn jhôrālen}.

\begin{cousinweb}
\textbf{\bn{ঝরানো}} (\emph{jhôrāno}) is the causative of \bn{ঝরা}, to fall away, which CDIAL lists for Bengali under a reconstructed form (CDIAL 5346): to make something fall away.
\end{cousinweb}

\begin{grammarlens}[title={What you've built}]
One more word for describing what you lived and felt.
\end{grammarlens}

\subsection*{Your turn}
\begin{itemize}
\raggedright
  \item \emph{Say it:} \emph{jhôrāno}
  \item \emph{Say it:} \emph{jhôrāno}, once more
  \item \emph{Say it:} \emph{bæyām kôre bābā pā̃ch keji ojôn jhôrālen.} — By exercising, my father shed five kilos.
  \item \emph{Recall:} say the Bengali for an addiction, then the Bengali for smoking, then say \emph{jhôrāno} again
\end{itemize}

\begin{tcolorbox}[breakable,colback=teal!4,colframe=teal!35!black,title={Before you move on}]
What does \bn{ঝরানো} mean? (\textquotedblleft{}To shed (weight, tears)\textquotedblright{}.) Say it once more.
\end{tcolorbox}
